\begin{definition}[The accessible two-site charge matrix]
    \label{def:peps_regular_charge_coordinates}
    \lean{TNLean.PEPS.regularChargeMatrix}
    \leanok
    Let $G$ be a finite group, $\chi:G\to\mathbb C$, and $p,x,y\in G$.
    Define the matrix $A_\chi(p;x,y)$ by
    \[
        A_\chi(p;x,y)_{r,s}
          =\begin{cases}
             \chi(px^{-1}r),&s=yx^{-1}r,\\
             0,&s\ne yx^{-1}r.
           \end{cases}
    \]
    This is the accessible-coordinate matrix in
    \cite[charge-detection proof]{Schuch2010PEPS}, lines 2470--2486.
    Here $p$ is the internal translation, and $x,y$ are the unknown vertex
    translations. No original-spin measurement is asserted.
\end{definition}

\begin{theorem}[Matrix units and the exact support-overlap pairing]
    \label{thm:peps_regular_charge_coordinates_pairing}
    \lean{TNLean.PEPS.regularChargeMatrix_eq_sum_single,
      TNLean.PEPS.regularChargeMatrix_inner}
    \leanok
    \uses{def:peps_regular_charge_coordinates}
    The accessible matrix is the literal matrix-unit sum
    \[
        A_\chi(p;x,y)
          =\sum_{g\in G}\chi(pg)|xg\rangle\langle yg|.
    \]
    For another function $\psi$ and translations $q,x',y'$, its exact
    Hilbert--Schmidt pairing is
    \begin{align}
      &\langle A_\chi(p;x,y),A_\psi(q;x',y')\rangle\\
      &\quad=\mathbf1_{yx^{-1}=y'(x')^{-1}}
        \sum_{r\in G}\overline{\chi(px^{-1}r)}\psi(q(x')^{-1}r).
    \end{align}
    Thus different relative translations have disjoint matrix support.
\end{theorem}
\begin{proof}\leanok
    At a row $r$, only $g=x^{-1}r$ contributes to the matrix-unit sum.
    The surviving column is $yx^{-1}r$. Different relative translations
    give distinct surviving columns at every row, so their support pairing
    vanishes. When they agree, sum the surviving entries.
\end{proof}

\begin{theorem}[Orthogonality and norm of the actual accessible charge matrices]
    \label{thm:peps_regular_charge_coordinates_orthogonality}
    \lean{TNLean.PEPS.regularChargeMatrix_orthogonal,
      TNLean.PEPS.regularChargeMatrix_normSq}
    \leanok
    \uses{thm:peps_regular_charge_coordinates_pairing,
      thm:peps_translated_character_inner_product}
    Let $\sigma,\tau$ be finite-dimensional irreducible complex representations,
    with $\sigma$ unitary. If $\chi_\sigma\ne\chi_\tau$, then for all
    $p,x,y,q,x',y'\in G$,
    \[
        \langle A_{\chi_\sigma}(p;x,y),
          A_{\chi_\tau}(q;x',y')\rangle=0.
    \]
    Every matrix of a unitary irreducible character has squared norm $|G|$,
    for every $p,x,y$. These are auxiliary accessible-coordinate identities
    in \cite[charge-detection proof]{Schuch2010PEPS}, lines 2470--2486;
    constructing the complete original-spin detector remains separate.
\end{theorem}
\begin{proof}\leanok
    For equal relative translations, apply translated character orthogonality
    with parameters $px^{-1}$ and $q(x')^{-1}$. For the norm,
    the same character appears on both sides and its translated squared norm
    is $|G|$.
\end{proof}

\begin{theorem}[The internal label is absorbed by unknown vertex translations]
    \label{thm:peps_regular_charge_coordinates_internal_label}
    \lean{TNLean.PEPS.regularChargeMatrix_internalLabel}
    \leanok
    \uses{def:peps_regular_charge_coordinates}
    For every function $\chi:G\to\mathbb C$ and every $p,q,x,y\in G$,
    \[
        A_\chi(q;xp^{-1}q,yp^{-1}q)=A_\chi(p;x,y).
    \]
    Thus changing the internal label can be absorbed by changing the two
    unknown vertex translations. This is the exact family-level identity
    underlying the final sentence of
    \cite[charge-detection proof]{Schuch2010PEPS}, lines 2483--2486.
    It is not a probabilistic or original-spin information-theoretic assertion.
\end{theorem}

\begin{proof}\leanok
    Cancel $q q^{-1}$ and $p^{-1}p$ in the relative column coordinate and
    character argument of every matrix entry.
\end{proof}

\begin{theorem}[Independent translation covariance and nonvanishing]
    \label{thm:peps_regular_charge_coordinates_covariance}
    \lean{TNLean.PEPS.regularChargeMatrix_translate_apply,
      TNLean.PEPS.regularChargeMatrix_ne_zero}
    \leanok
    \uses{thm:peps_regular_charge_coordinates_orthogonality,
      def:peps_regular_charge_coordinates}
    For every coefficient function $\chi$ and every $p,x,y,z,w,r,s\in G$,
    \[
        A_\chi(p;zx,wy)_{r,s}
          =A_\chi(p;x,y)_{z^{-1}r,w^{-1}s}.
    \]
    Thus the two unknown translations act by their actual row and column
    permutations. If $\chi$ is a unitary irreducible character, then every
    matrix $A_\chi(p;x,y)$ is nonzero. These are auxiliary identities from
    the accessible charge coordinates of
    \cite[charge detection]{Schuch2010PEPS}, lines 2470--2486.
\end{theorem}
\begin{proof}\leanok
    Substitute the translated row and column coordinates and cancel the
    inverse translations. For nonvanishing, the squared norm is $|G|>0$.
\end{proof}
